% Para usar o diagrama abaixo no artigo, adicione no preambulo:
% \usepackage{tikz}
% \usetikzlibrary{arrows.meta,positioning,fit,shapes.geometric,backgrounds}

\section{ARQUITETURA}

A arquitetura adotada foi organizada como um fluxo modular de consulta documental baseado em Retrieval-Augmented Generation (RAG). O objetivo principal da solução é responder perguntas em linguagem natural a partir de um corpus fechado de documentos doutrinários e normativos, preservando a rastreabilidade das fontes e evitando que o modelo de linguagem atue como fonte autônoma de informação. Por essa razão, o desenho arquitetural separa explicitamente a preparação do corpus, a recuperação de evidências, a geração condicionada ao contexto e a apresentação da resposta ao usuário.

A Figura~\ref{fig:arquitetura-modular} apresenta uma visão geral da arquitetura. O fluxo online começa na interface web, passa pela API da aplicação e aciona o pipeline RAG responsável por normalizar a pergunta, recuperar trechos relevantes, aplicar reordenação e verificar se existe evidência documental suficiente para sustentar a resposta. Em paralelo, o fluxo offline prepara os documentos do corpus, gerando chunks estruturais, metadados e índices de busca que são consumidos no momento da consulta.

\begin{figure*}[t]
\centering
\begin{tikzpicture}[
    font=\scriptsize,
    box/.style={draw, rounded corners=2pt, align=center, minimum width=2.7cm, minimum height=0.82cm, fill=white},
    group/.style={draw, rounded corners=2pt, inner sep=7pt, fill=#1},
    arrow/.style={-{Latex[length=2mm]}, thick},
    data/.style={draw, cylinder, shape border rotate=90, aspect=0.23, align=center, minimum height=1.1cm, minimum width=1.45cm, fill=green!8}
]

\node[box, fill=blue!8] (user) {Usuário\\pergunta em linguagem natural};
\node[box, fill=blue!8, right=0.9cm of user] (front) {Frontend web\\React + Vite};
\node[box, fill=purple!10, right=0.9cm of front] (api) {API backend\\FastAPI + Uvicorn};
\node[box, fill=purple!10, right=0.9cm of api] (rag) {Orquestração RAG\\Python};
\node[box, fill=yellow!15, right=0.9cm of rag] (answer) {Resposta\\com fontes ou recusa};

\node[box, fill=orange!12, below=1.15cm of rag, xshift=-3.25cm] (query) {Normalização\\da pergunta};
\node[box, fill=orange!12, right=0.75cm of query] (retrieval) {Pipeline de\\recuperação};
\node[box, fill=orange!12, right=0.75cm of retrieval] (evidence) {Política de\\evidência};
\node[box, fill=orange!12, right=0.75cm of evidence] (prompt) {Prompt + citações\\contexto controlado};
\node[box, fill=orange!12, right=0.75cm of prompt] (llm) {Modelo de linguagem\\OpenAI API};

\node[data, below=1.2cm of retrieval, xshift=-1.35cm] (chroma) {ChromaDB\\índice vetorial};
\node[data, right=0.65cm of chroma] (bm25) {BM25\\índice lexical};
\node[box, fill=green!10, right=0.65cm of bm25] (rerank) {Fusão híbrida\\RRF + reranking};
\node[box, fill=green!10, right=0.65cm of rerank] (hier) {Recuperação\\hierárquica};

\node[box, fill=gray!10, below=1.25cm of chroma, xshift=-1.4cm] (pdfs) {PDFs do corpus\\controlado};
\node[box, fill=gray!10, right=0.55cm of pdfs] (extract) {Extração\\PyMuPDF};
\node[box, fill=gray!10, right=0.55cm of extract] (chunk) {Normalização e\\chunking estrutural};
\node[box, fill=gray!10, right=0.55cm of chunk] (meta) {Chunks JSONL\\metadados e fontes};
\node[box, fill=gray!10, right=0.55cm of meta] (embed) {Embeddings\\e índices};

\node[box, fill=red!8, below=1.2cm of prompt] (audit) {Relatórios e auditoria\\recuperação, geração e fontes};
\node[box, fill=red!8, right=0.65cm of audit] (deploy) {Infraestrutura\\Docker + ambiente de execução};

\draw[arrow] (user) -- (front);
\draw[arrow] (front) -- node[above]{HTTP/JSON} (api);
\draw[arrow] (api) -- (rag);
\draw[arrow] (rag) -- (answer);
\draw[arrow] (answer) .. controls +(0,1.2) and +(0,1.2) .. (front);

\draw[arrow] (rag) -- (query);
\draw[arrow] (query) -- (retrieval);
\draw[arrow] (retrieval) -- (evidence);
\draw[arrow] (evidence) -- (prompt);
\draw[arrow] (prompt) -- (llm);
\draw[arrow] (llm) -- (answer);

\draw[arrow] (retrieval) -- (chroma);
\draw[arrow] (retrieval) -- (bm25);
\draw[arrow] (chroma) -- (rerank);
\draw[arrow] (bm25) -- (rerank);
\draw[arrow] (rerank) -- (hier);
\draw[arrow] (hier) -- (evidence);

\draw[arrow] (pdfs) -- (extract);
\draw[arrow] (extract) -- (chunk);
\draw[arrow] (chunk) -- (meta);
\draw[arrow] (meta) -- (embed);
\draw[arrow] (embed) -- (chroma);
\draw[arrow] (embed) -- (bm25);

\draw[arrow] (rag) -- (audit);
\draw[arrow] (api) |- (deploy);

\begin{scope}[on background layer]
\node[group=blue!3, fit=(user)(front)] (gfront) {};
\node[group=purple!4, fit=(api)(rag)(answer)(query)(retrieval)(evidence)(prompt)(llm)] (gback) {};
\node[group=green!3, fit=(chroma)(bm25)(rerank)(hier)] (gret) {};
\node[group=gray!4, fit=(pdfs)(extract)(chunk)(meta)(embed)] (goff) {};
\node[group=red!3, fit=(audit)(deploy)] (gops) {};
\end{scope}

\node[above=0.05cm of gfront] {Camada de interação};
\node[above=0.05cm of gback] {Backend e orquestração};
\node[above=0.05cm of gret] {Recuperação e ranqueamento};
\node[above=0.05cm of goff] {Preparação offline do corpus};
\node[above=0.05cm of gops] {Operação e rastreabilidade};

\end{tikzpicture}
\caption{Visão geral da arquitetura modular do sistema proposto.}
\label{fig:arquitetura-modular}
\end{figure*}

\subsection{Visão geral da solução}

O sistema foi estruturado em três camadas principais: interface, backend e pipeline RAG. A interface web permite que o usuário formule perguntas e visualize a resposta acompanhada de fontes. O backend, implementado como uma API, centraliza os endpoints de consulta e encaminha as requisições para a camada de geração. Essa camada não envia a pergunta diretamente ao modelo de linguagem; antes disso, executa etapas de recuperação, filtragem e validação de evidências.

Essa separação reduz o acoplamento entre a experiência conversacional e os componentes de inteligência artificial. Assim, mudanças no mecanismo de recuperação, no modelo de linguagem ou na política de citação podem ser feitas sem alterar a lógica central da interface. Além disso, a arquitetura favorece a manutenção do corpus como uma base controlada, evitando uploads livres ou fontes externas não verificadas durante a execução.

\subsection{Preparação e indexação do corpus}

A preparação do corpus ocorre em um fluxo offline. Os documentos em PDF são submetidos à extração textual, normalização e segmentação em chunks. Diferentemente de uma divisão puramente baseada em tamanho, a segmentação considera a estrutura dos documentos, como capítulos, seções, artigos, perguntas e respostas. Essa decisão é importante porque o significado de um trecho doutrinário ou normativo frequentemente depende da unidade documental em que ele aparece.

Após a segmentação, cada chunk recebe metadados associados à sua origem, como documento, tipo de corpus, localização textual e referência de fonte. Esses metadados permitem rastrear a proveniência das respostas e também aplicar restrições durante a recuperação. Em seguida, os chunks são persistidos em arquivos estruturados e utilizados para alimentar dois índices complementares: um índice vetorial, voltado à recuperação semântica, e um índice lexical, voltado à correspondência de termos específicos.

\subsection{Recuperação e ranqueamento}

No momento da consulta, a pergunta do usuário passa por uma etapa de normalização leve e é encaminhada ao pipeline de recuperação. A arquitetura combina busca vetorial, baseada em embeddings, e busca lexical, baseada em BM25. A busca vetorial contribui para encontrar trechos semanticamente relacionados à pergunta, mesmo quando há diferença de vocabulário entre a consulta e o documento. A busca lexical, por sua vez, preserva a importância de termos técnicos, expressões doutrinárias e categorias normativas que não devem ser diluídas apenas por aproximação semântica.

Os resultados das duas estratégias são combinados por fusão de rankings e, em seguida, reordenados por um modelo de reranking. Essa etapa busca priorizar os trechos que respondem mais diretamente à pergunta formulada. Depois do ranqueamento, a recuperação hierárquica expande os trechos selecionados para unidades documentais mais amplas quando necessário, preservando contexto suficiente para a geração da resposta.

\subsection{Geração condicionada a evidências}

A geração da resposta é condicionada aos contextos recuperados. Antes de construir o prompt final, uma política de evidência verifica se os trechos encontrados oferecem suporte suficiente para responder à pergunta. Quando a evidência é considerada adequada, o prompt inclui a pergunta do usuário, os trechos selecionados e instruções para produzir uma resposta vinculada às fontes. Quando a evidência é insuficiente, o sistema deve sinalizar a limitação em vez de produzir uma resposta especulativa.

Essa política é central para o domínio do trabalho, pois impede que o modelo de linguagem seja tratado como autoridade independente. O modelo atua como componente de síntese textual: ele organiza e explica as evidências recuperadas, mas a validade da resposta depende dos documentos consultados. Ao final, as fontes utilizadas são formatadas e retornadas junto com a resposta, permitindo que o usuário verifique a origem da informação apresentada.

\subsection{Implantação e rastreabilidade}

A aplicação foi organizada para execução reprodutível por meio de empacotamento com Docker. O backend, a interface web e os artefatos necessários ao funcionamento do corpus podem ser executados como uma unidade de aplicação, reduzindo diferenças entre ambiente local e ambiente de implantação. As variáveis de ambiente controlam chaves de API, caminhos do corpus, diretórios do índice vetorial e parâmetros do pipeline.

Além da execução, a arquitetura prevê relatórios e registros associados às etapas de preparação do corpus, recuperação e geração. Esses artefatos apoiam a auditoria do comportamento do sistema, pois permitem investigar quais documentos foram processados, quais trechos foram recuperados e quais fontes sustentaram determinada resposta. Dessa forma, a solução não se limita a entregar uma resposta conversacional, mas mantém mecanismos de rastreabilidade compatíveis com um domínio documental sensível.
